% cap02/2.6_protocolos.tex
\section{Evaluación Topológica del Tráfico de Red: UDP frente a la Integridad Estricta de TCP}

En la búsqueda imperativa de la excelencia arquitectónica para salvaguardar criptográfica e infaliblemente los transaccionales activos métricos del tercer sector, el ingeniero de software se enfrenta ineludiblemente al dilema clásico de la capa de transporte del modelo OSI: la elección diametral entre el protocolo ligero,, veloz y ciego UDP (\textit{User Datagram Protocol}) frente a la, estricta, síncrona y algorítmicamente exhaustiva fidelidad métrica y relacional del protocolo TCP (\textit{Transmission Control Protocol}). El presente apartado no busca enumerar trivialmente sus rígidas estructuras de cabecera; por el contrario, se adentra matemática y en disectar cómo la naturaleza eléctrica de estos protocolos colisiona y brutalmente con los dogmáticos requisitos de persistencia segura, y (\textit{ACID}) exigidos por la plataforma e interactiva de logística civil denominada \textit{Democra}.

\subsection{El Espejismo de la Velocidad UDP}

El protocolo \textit{UDP} es reverenciado en ecosistemas de transmisión continua (como el masivo flujo métrico de streaming de video o la latencia crítica de radares militares de telemetría cíclica). Su brutal velocidad de ruteo radica y arquitectónicamente en una osada, deliberada y carencia de mecanismos de control. Al transmitir ciegamente un paquete, UDP jamás establece una conexión formal (\textit{connectionless}); llanamente y algorítmicamente dispara y rutea el datagrama hacia la ciega inmensidad de la red, sin importarle matemática, o si el receptor logró capturarlo, si el paquete llegó corrupto o si los bytes se perdieron irrevocablemente y para siempre en el caótico ruteo de un switch saturado.

En la y rigurosa teoría militar de DDS, esta pérdida es tolerada porque el siguiente paquete de radar reescribirá ciegamente la posición en milisegundos. No obstante, en la y transaccional de gestión de bases de donativos solidarios, esta ceguera de UDP es letal y inaceptable. Si un mensaje UDP conteniendo la instrucción de descontar cien kilos de arroz es dropeado en la red rural, la base de datos relacional \textit{Multi-Tenant} jamás lo registrará, creando instantáneamente un letal descuadre de inventario. La abrumadora métrica no permite aproximaciones; cada variable de \textit{stock} debe ser matemáticamente perfecta.

\subsection{La Fortaleza Inquebrantable de TCP}

Ante la innegable y letal paradoja falibilidad de UDP, la majestuosa arquitectura de la aplicación y \textit{Democra} abraza y férreamente el estándar \textit{Transmission Control Protocol} (TCP). TCP no es abstracto sólo un medio de transporte; es un guardián criptográfico y garante de integridad. 

Antes de enviar un sólo byte, el implacable y protocolo TCP ejecuta el famoso \textit{Three-Way Handshake}, un saludo a tres vías que asegura que tanto el cliente \textit{Frontend} como el servidor \textit{Backend} estén listos. Durante la transmisión, tcp añade un código de secuencia a cada byte. Si la aplicación envía la orden de descontar arroz, el servidor debe devolver un acuse de recibo (\textit{ACK}). Si este \textit{ACK} no llega, el cliente TCP retransmite de forma silenciosa y segura el mensaje hasta confirmar su entrega.

Esta majestuosa y fidelidad TCP es la única base y cimiento capaz de soportar las y restricciones de \textit{Row-Level Security} (RLS) en la base de datos. Sólo un protocolo TCP puede transportar los densos y masivos \textit{JSON Web Tokens} (JWT), garantizando que el criptográfico identificador del inquilino de la ONG llegue intacto e inalterable al motor SQL, previniendo matemáticamente que la exfiltración de datos ocurra por caída de bytes. Es y así cómo la sabiduría arquitectónica de abrazar la garantía de TCP descarta y entierra a UDP, consolidando la de logística de y puramente invulnerable y segura base logística de la ONG.
